\begin{table}[htbp]\centering
\caption{The espoused--enacted gap and Progress~8}
\label{tab:gaps}
\def\sym#1{\ifmmode^{#1}\else\(^{#1}\)\fi}
\small

\footnotesize\setlength{\tabcolsep}{4pt}
\begin{tabular}{lccc}
\toprule
 & Signed gap & Absolute gap & \shortstack{Absolute gap\\(+ component quadratics)} \\
\midrule
\multicolumn{4}{l}{\textit{Visited sample (visit vs interview, $n=99$)}} \\
\addlinespace[2pt]
Warmth & 0.047 & 0.032 & 0.011 \\
 & (0.040) & (0.046) & (0.058) \\
Strictness & 0.065 & 0.002 & $-$0.015 \\
 & (0.050) & (0.061) & (0.061) \\
\bottomrule
\end{tabular}
\begin{minipage}{\linewidth}
\smallskip
\footnotesize\textit{Notes:} Each row regresses average Progress 8 on the enacted (observed) score, the espoused-minus-enacted gap, and the primary control set, on the visited schools. The first column enters the signed gap, the second the absolute gap, and the third re-enters the absolute gap alongside squared terms in both the espoused and the enacted score. An absolute gap could predict progress simply because progress is a non-linear function of one of the scores. The squared terms remove that explanation. Standard errors in parentheses (HC3).
\sym{*} \(p<0.10\), \sym{**} \(p<0.05\), \sym{***} \(p<0.01\).
Scores are z-standardised before the gap is formed; late-entry schools excluded. 
\end{minipage}
\end{table}
